\documentclass[11pt,a4paper]{article}
\usepackage[margin=23mm]{geometry}
\usepackage[T1]{fontenc}
\usepackage{lmodern,microtype,amsmath,amssymb,amsthm,mathtools,booktabs,array,xcolor}
\usepackage[colorlinks=true,linkcolor=blue!45!black,urlcolor=blue!45!black]{hyperref}
\newtheorem{theorem}{Theorem}[section]
\newtheorem{lemma}[theorem]{Lemma}
\newtheorem{proposition}[theorem]{Proposition}
\newtheorem{corollary}[theorem]{Corollary}
\theoremstyle{remark}\newtheorem{remark}[theorem]{Remark}
\newcommand{\R}{\mathbb R}\newcommand{\C}{\mathbb C}
\newcommand{\ii}{\mathrm i}\newcommand{\dd}{\,\mathrm d}
\newcommand{\norm}[1]{\left\|#1\right\|}
\newcommand{\ip}[2]{\left\langle #1,#2\right\rangle}
\newcommand{\Rea}{\operatorname{Re}}\newcommand{\Ima}{\operatorname{Im}}
\newcommand{\dist}{\operatorname{dist}}\newcommand{\spec}{\operatorname{spec}}
\newcommand{\Var}{\operatorname{Var}}\newcommand{\rank}{\operatorname{rank}}
\setlength{\parskip}{3pt}
\title{Zero-momentum global dynamics and three-dimensional traveling waves\\for the mass-critical biharmonic NLS}
\author{Shuxing Yang}
\date{7 October 2026}
\begin{document}
\maketitle
\begin{abstract}
This continuation develops three results for the focusing mass-critical biharmonic nonlinear Schr\"odinger equation. First, a finite-rank, spatially localizable modification of the linearized energy permits outgoing excited-wave gluing with arbitrarily small nonzero separation speed. Applied to the previously certified four-dimensional negative-energy branch, it gives bounded forward-global solutions with exactly zero momentum, finite variance, and negative energy. Second, whole-space interval estimates complete the three-dimensional profile, negative inverse pairing, and all symmetry-reduced angular channels; an implicit-function argument yields negative-energy traveling waves in dimension three. Third, a purely analytic interaction-virial and mass-transport argument replaces a uniform kinetic bound by its first time average and forces separation on entire time intervals. Uniform mass localization instead forces quantitative averaged kinetic grow-up. All numerical conclusions refer to explicit frozen data, not to plots or finite-box eigenvalues alone.
\end{abstract}

\noindent\textbf{Profile inputs.} The four-dimensional application retains the explicit companion profile certificates as an input; the three-dimensional certificates are completed here.  A conflict with a published cylindrical blowup statement is addressed in Section~\ref{sec:context}; identifying a gap in that proof alone would not establish our construction.

\section{Equation, conventions, and principal statements}
We use
\begin{equation}\label{eq:pde}
 \ii u_t=\Delta^2u-|u|^{8/N}u,\qquad u(t)\in H^2(\R^N),
\end{equation}
and the conserved quantities
\[
 M(u)=\int|u|^2,\quad
 E(u)=\frac12\norm{\Delta u}_2^2-
 \frac1{2+8/N}\norm u_{2+8/N}^{2+8/N},\quad
 P(u)=\Ima\int\bar u\nabla u.
\]
Put $A=\Delta^2+1$ and $\norm h_Z^2=\norm{\Delta h}_2^2+\norm h_2^2$.
A wave $u(t,x)=e^{\ii t}U_v(x-vte_1)$ solves~\eqref{eq:pde} precisely when
\begin{equation}\label{eq:profile}
 A U_v+\ii v\partial_1U_v-|U_v|^{8/N}U_v=0.
\end{equation}
For an admissibly localized mass-critical profile, variation under
$U(x)\mapsto\lambda^{N/2}U(\lambda x)$ in
$E+M/2-vP_1/2$ gives
\begin{equation}\label{eq:pohozaev}
 E(U_v)=\frac{vP_1(U_v)}8.
\end{equation}
There is no Galilean reduction that forces $vP_1$ to be positive for this fourth-order equation.

\begin{theorem}[Four-dimensional zero-momentum construction]\label{thm:pair}
Let $Q_v$ be the previously certified smooth exponentially localized negative-energy cubic BNLS profile, for any fixed admissible $0<v\le10^{-6}$. There exist $T,\gamma>0$ and an exact solution such that, with $Q_{-v}(x)=Q_v(-x)$,
\[
 R(t,x)=e^{\ii t}\bigl(Q_v(x-vte_1)+Q_{-v}(x+vte_1)\bigr),
\]
\begin{equation}\label{eq:weightedpair}
 \norm{u(t)-R(t)}_{H^{10}}+
 \sum_{j=1}^4\norm{x_j(u(t)-R(t))}_{H^6}\le Ce^{-\gamma t}
 \qquad(t\ge T).
\end{equation}
One can take $u$ inversion-even and invariant under transverse $O(3)$. It has
\[
 P(u)=0,\qquad E(u)=2E(Q_v)<0,\qquad M(u)=2M(Q_v),
 \qquad\sup_{t\ge T}\norm{u(t)}_{H^{10}}<\infty.
\]
Its variance is finite at every finite time and satisfies
\begin{equation}\label{eq:variance}
 \int|x|^2|u(t,x)|^2\dd x
 =2M(Q_v)v^2t^2+2\int|x|^2|Q_v(x)|^2\dd x+O(te^{-\gamma t}).
\end{equation}
Time translation gives a forward-global solution with initial time zero.
\end{theorem}

The profile input already proves $E(Q_v)<-18v^2$ on the displayed explicit speed interval. The gluing theorem itself needs neither an exact Morse count nor the final radial Hamiltonian quartet. Starting separation and all constants may deteriorate as $v\downarrow0$.

\begin{theorem}[Three-dimensional negative-energy waves]\label{thm:n3}
The frozen three-dimensional rational profile $q$ specified in Section~\ref{sec:n3} has an exact real radial solution $Q$ with
\[
 A Q-|Q|^{8/3}Q=0,\qquad\norm{Q-q}_Z<4\times10^{-6}.
\]
On $X=\{h\in H^2(\R^3;\C):h(-x)=\overline{h(x)}\}$, its real linearization is invertible with norm less than $200$ from $X^*$ to $X$. Moreover,
\begin{equation}\label{eq:n3I}
 I=\ip{\partial_1Q}{(A-|Q|^{8/3})^{-1}\partial_1Q}<-1280.
\end{equation}
There is a $C^3$ branch solving~\eqref{eq:profile} for $|v|\le10^{-9}$ with
\[
 U_v=Q-\ii v(A-|Q|^{8/3})^{-1}\partial_1Q+o_Z(v),
 \qquad P_1(U_v)=2Iv+o(v).
\]
Quantitatively, $\norm{U_v-Q}_Z\le1000|v|$ and $E(U_v)<-180v^2$ for $0<|v|\le10^{-9}$. Thus $e^{\ii t}U_v(x-vte_1)$ is an exact global $H^2$ solution with constant $H^2$ norm.
\end{theorem}

Thus the proposed unrestricted implication $E<0\Rightarrow$ finite-time collapse has counterexamples in both dimensions three and four, subject to the supplied computer-assisted proof chain. The four-dimensional conclusion survives adding zero momentum and finite variance. This does not classify negative-energy data or settle every fixed dimension separately.

\section{A compact finite-rank energy correction}
This is the central analytic mechanism in Theorem~\ref{thm:pair}. Work over real Hilbert spaces and fix a smooth exponentially localized cubic profile in $\R^4$ with $0<|v|<1/2$. Let
\[
 L_v=A+\ii v\partial_1-\mathcal V_v,\quad
 \mathcal V_vh=2|Q_v|^2h+Q_v^2\bar h,\quad
 J=-\ii,\quad B_v=JL_v.
\]
The free symbol $|\xi|^4+1-v\xi_1$ has minimum
$1-3(|v|/4)^{4/3}>0$. The decaying potential is compact from $H^2$ to $H^{-2}$, so $L_v$ has finite negative index, finite kernel, and a coercive positive complement. The graph domains are $D(B_v^k)=H^{2+4k}$. These facts verify the hypotheses for the structural decomposition of Lin--Zeng~\cite{LZ}.

\begin{lemma}[Localizable almost-conserved metric]\label{lem:metric}
For every $\varepsilon>0$ there is a self-adjoint finite-rank form $K_\varepsilon$ with smooth compactly supported kernels such that
\[
 c_\varepsilon\norm h_{H^2}^2\le
 q_\varepsilon(h):=\ip{(L_v+K_\varepsilon)h}{h}
 \le C_\varepsilon\norm h_{H^2}^2.
\]
On the invariant center space $E^c$ of $B_v$,
\begin{equation}\label{eq:metricder}
 \left|\ip{(K_\varepsilon B_v+B_v^*K_\varepsilon)h}{h}\right|
 \le\varepsilon q_\varepsilon(h).
\end{equation}
The inequality is a bounded quadratic-form statement on $E^c\subset H^2$.
\end{lemma}

\begin{proof}
Use the full triangular decomposition $H^2=\bigoplus_{j=0}^6X_j$ from \cite[Theorem 2.1]{LZ}. Here $X_0=\ker L_v$, $X_5=E^u$, $X_6=E^s$, all blocks except $X_3$ are finite-dimensional and lie in every graph domain. On the center $X_0\oplus\cdots\oplus X_4$, the only triangular arrows are
\[
 4\longrightarrow\{2,3,1,0\},\qquad
 \{2,3\}\longrightarrow\{1,0\},\qquad1\longrightarrow0.
\]
The diagonal operators on blocks $2,3$ are skew in the coercive metrics $-L_2,L_3$; those on $1,4$ are finite matrices with imaginary spectrum. Rescaling their Jordan chains gives positive metrics $G_1,G_4$ with arbitrarily small symmetric generator parts. Set
\[
 q_\delta=a_0|x_0|^2+a_1G_1[x_1,x_1]-L_2[x_2,x_2]
              +L_3[x_3,x_3]+a_4G_4[x_4,x_4].
\]
Take $a_4$ large, then $a_1$ small, then $a_0$ smaller. Each off-diagonal norm in these coordinates is its original norm times the square root of the target/source weight ratio. Thus $|q_\delta'|\le\delta q_\delta$. The infinite-dimensional metric remains exactly $L_3$. Extending positively on the two hyperbolic blocks changes $L_v$ by a finite-rank form only.

The dual functions of this correction are smooth. Coordinates on $X_1,X_4$ are recovered by pairing against $L_vX_4,L_vX_1$, respectively; coordinates on $X_2$ use $L_vX_2$, and hyperbolic coordinates use the opposite hyperbolic block. For the kernel, use
\[
 P_0=\operatorname{Orth}^{H^2}_{X_0}(I-P_5-P_6).
\]
This accounts for the fact that kernel orthogonality need not extend to hyperbolic blocks. All these duals belong to every Sobolev space, so the correction is a finite sum
$\sum c_{jk}\ip{f_j}{h}\ip{f_k}{h}$ with smooth Sobolev $f_j$.

Replace $f_j$ by smooth compact cutoffs $\chi_Rf_j$. They converge in each fixed Sobolev norm. Both the form error and its generator-commutator error tend to zero: in a rank-one term move $B_v$ to the smooth dual function, using $B_v^*:H^2\to H^{-2}$. First fix the small-rate metric, then choose the support radius relative to its coercivity. Positivity and~\eqref{eq:metricder} follow. No decay assumption on embedded imaginary eigenmodes is used.
\end{proof}

For the existing branch the correction rank can be bounded by
$\dim\ker L_v+2n^-(L_v)=5+2\cdot36=77$. The argument needs only finiteness; this number is not a hypothesis for gluing.

\section{Moving-center gluing and weighted regularity}
\subsection{The hyperbolic projection and center evolution}
Only the finitely many off-axis generalized eigenvectors and their duals must be transported exactly. They are smooth and exponentially localized: the free fourth-order resolvent is invertible off the imaginary spectrum, small exponential conjugation is a relative graph perturbation, and exterior potential terms can be absorbed. Induction treats Jordan chains.

Translate synthesis and analysis maps to the two centers and correct the exponentially small off-diagonal Gram matrix. Write them $Z(t),W(t)$ with $WZ=I$, and let $B_0$ be the direct sum of the two hyperbolic generators. With
\[
 S=Q_v(\cdot-vte_1)+Q_{-v}(\cdot+vte_1),\quad
 \mathcal A=J(A-\mathcal V_S),
\]
the errors in
\[
 \dot Z+ZB_0-\mathcal AZ=E_Z,\qquad
 \dot W+W\mathcal A-B_0W=E_W
\]
are exponentially small in every fixed graph norm used here. Define
\[
 \Pi=I-ZW,\qquad\mathcal C=\mathcal A-ZE_W.
\]
Then $w_t=\mathcal Cw$ preserves $W(t)w(t)=0$.

Choose a real axial velocity field $b$ equal to $v$ near the right core and $-v$ near the left, with transition length comparable to $vt$. Set
\[
 D_b=b\partial_1+\tfrac12\partial_1b,\qquad
 \mathcal H_\varepsilon=A-\mathcal V_S+\ii D_b
                         +K_\varepsilon^+(t)+K_\varepsilon^-(t).
\]
Its coefficients satisfy $\norm{\partial_x^jb}_\infty=O_v(t^{-j})$ and $\norm{b_t}_\infty=O_v(t^{-1})$. A square partition $\chi_+^2+\chi_-^2+\chi_0^2=1$ has core cutoffs equal to one on the compact correction supports. Fourth-order IMS and the free exterior positivity give, for sufficiently large $t$,
\begin{align}
 Q_t(w):=\ip{\mathcal H_\varepsilon w}{w}&\asymp_\varepsilon\norm w_{H^2}^2,\label{eq:ims1}\\
 q_\varepsilon^+(\chi_+w)+q_\varepsilon^-(\chi_-w)&\le2Q_t(w).\label{eq:ims2}
\end{align}
The constant $2$ in~\eqref{eq:ims2}, after increasing the starting time, is essential: replacing local energies by a norm would insert an uncontrolled condition number in front of $\varepsilon$.

For $\mathcal H_0=A-\mathcal V_S+\ii D_b$, direct real-form differentiation gives
\begin{equation}\label{eq:comm}
 \dot{\mathcal H}_0+\mathcal H_0\mathcal A+\mathcal A^*\mathcal H_0
 =\ii D_{b_t}+[D_b,\Delta^2]-\bigl((\mathcal V_S)_t+b\partial_1\mathcal V_S\bigr).
\end{equation}
Transport cancels the last term in the cores; tails are exponentially small. Integrating the commutator as a quadratic form bounds it by $C_vt^{-1}\norm w_{H^2}^2$. The translation derivative of each $K_\varepsilon^\pm$ cancels the matching transport. The remaining compact form is exactly the single-wave form in~\eqref{eq:metricder}, up to exponential errors.

If $Ww=0$, the local vectors $\chi_\pm w$ have only exponentially small single-wave hyperbolic coordinates, since those duals decay exponentially. Project them onto the corresponding centers and use~\eqref{eq:ims2}. Consequently
\[
 |Q_t'(w)|\le(2\varepsilon+C_\varepsilon/t+C_\varepsilon e^{-\kappa t})Q_t(w),
\]
and the center propagator satisfies
\begin{equation}\label{eq:center}
 \norm{U_C(t,s)}_{H^2\to H^2}
 \le C_\varepsilon e^{\varepsilon|t-s|}
       \max(t/s,s/t)^{m_\varepsilon}.
\end{equation}
This provides any prescribed small exponential rate, even for arbitrarily small fixed $v$.

\subsection{Graph norms and one spatial moment}
Let $\mathcal D=\Pi\mathcal A\Pi$. The difference $\mathcal D-\mathcal A$ is finite rank and smoothing. On the moving center,
\[
 \norm w_{H^6}\asymp\norm w_{H^2}+\norm{\mathcal Dw}_{H^2}.
\]
For $z=\mathcal Dw$,
\[
 z_t=\mathcal Cz+\mathcal Fw,\qquad
 \mathcal F=\dot{\mathcal D}+[\mathcal D,\mathcal C].
\]
Here $\mathcal F$ is an order-zero smooth multiplier plus smooth finite-rank terms. It is bounded on each fixed Sobolev space, and $\mathcal Fw$ belongs to the moving center. Duhamel gives~\eqref{eq:center} in $H^6$, with an additional polynomial factor. Iterating
\[
 (\mathcal D^k)_t+[\mathcal D^k,\mathcal C]
 =\sum_{j=0}^{k-1}\mathcal D^j\mathcal F\mathcal D^{k-1-j}
\]
gives the same small exponential rate in $H^{2+4k}$.

For position weights set $z_j=\Pi x_jw$ and
$B=\dot\Pi+\Pi\mathcal C-\mathcal C\Pi$. Since $B\Pi=0$,
\begin{equation}\label{eq:position}
 (z_j)_t-\mathcal Cz_j=B(I-\Pi)x_jw+\Pi[x_j,\mathcal C]w.
\end{equation}
The raw hyperbolic constraints vanish individually. A translated dual therefore satisfies
\[
 \ip{\psi(\cdot-c_k(t))}{x_jw}
 =\ip{(x_j-c_{k,j}(t))\psi(\cdot-c_k(t))}{w}.
\]
This removes the growing center displacement. The right side of~\eqref{eq:position} is bounded from $H^{10}$ to $H^6$, uniformly for large $t$: the free position commutator has order three, same-core finite-rank centers cancel, and cross-core polynomial factors multiply exponential errors. Thus the center bound extends to
\[
 Y=H^{10}\cap\{h:x_jh\in H^6,\ 1\le j\le4\},
\]
again with only additional polynomial factors. On finite hyperbolic coordinates, the $Y$ norm is comparable to $(1+t)|a|$, so transport ratios also cost only a polynomial in $|t-s|$.

\subsection{The exact nonlinear solution}
The interaction defect is $O_Y(e^{-\kappa t})$, after decreasing $\kappa$ to absorb position factors. Let $\alpha>0$ be a hyperbolic decay rate, omitting it if the hyperbolic space is empty. Choose
\[
 0<\gamma<\min(\kappa,\alpha),\qquad 0<\varepsilon<\gamma/2,
\]
and only then choose the compact metric and a sufficiently large $T$. In the exponentially weighted norm $\sup_{t\ge T}e^{\gamma(t-T)}\norm h_Y$, the center integral backward from infinity converges because
\[
 \int_0^\infty(1+r)^{m_\varepsilon}e^{-(\gamma-\varepsilon)r}\dd r<\infty.
\]
Stable coordinates are integrated forward and unstable coordinates backward. The remaining exponential couplings are absorbed by a Neumann series. The cubic nonlinear remainder is quadratic and Lipschitz on $Y$; each position factor can be assigned to a correction factor, for example
\[
 \norm{x_j(S|h|^2)}_{H^6}
 \le C\norm S_{H^6}\norm h_{H^{10}}\norm{x_jh}_{H^6}.
\]
The uniform unweighted norms of $S$ suffice. A contraction for sufficiently large $T$ proves~\eqref{eq:weightedpair}.

Reflection can be imposed throughout. Compact transverse symmetry can be imposed by averaging the bounded Green right inverse over $O(3)$ and restricting the contraction to invariant functions. This averages the right inverse, not the finite-rank metric. Conservation and asymptotic additivity yield the exact $M,E,P$ identities. The profile density is inversion-even, with zero intrinsic centroid; weighted convergence proves~\eqref{eq:variance}.

For sufficiently small $v$, the one-wave stable dimension is exactly $31$, using the additional certificate in Section~\ref{sec:quartet}. Prescribing small reflected stable coordinates at $T$ gives an injective family with $31$ real parameters and identical $M,E,P$. This family retains inversion symmetry and finite variance. Its full parameter count is not asserted after additionally imposing transverse $O(3)$ symmetry, and it is not an open set of bounded data.

\section{Whole-space three-dimensional certification}\label{sec:n3}
The reference file is \texttt{round4\_n3\_mixed\_exact.json}, with SHA-256
\begin{center}\small\texttt{5d2a86dc991fe960e7036ca651634c556ded9da4aa7d1ae357ab5f4f656d526b}.
\end{center}
It specifies rational compactly supported radial polynomials on $[0,54]$. Mixed auxiliary fields reduce fourth-order residuals to $L^2$ polynomial integrals. The existing verified bounds are
\[
 \norm{Aq-|q|^{8/3}q}_{Z^*}<1.3\times10^{-7},\qquad
 \norm q_Z<137,\qquad\norm q_\infty<6.3.
\]
The embedding $\norm h_\infty\le(4\sqrt2\pi)^{-1/2}\norm h_Z<\norm h_Z/4$ is used below.

\subsection{Radial inverse and exact nonlinear profile}
Let $V=(11/3)|q|^{8/3}$. Under the unitary radial map $f\mapsto\sqrt{4\pi}\,rf$, $A$ becomes $H^2+1$ on the half-line, where $H=-\partial_r^2$ is Dirichlet at zero. Compare
$\sqrt V(H^2+1)^{-1}\sqrt V$ with its Dirichlet-box resolvent at $R=64$.
For $\kappa=(1+\ii)/\sqrt2$, the second-order kernel difference is the rank-one expression
\[
 \frac{e^{-\kappa R}}{\kappa\sinh(\kappa R)}
       \sinh(\kappa r)\sinh(\kappa s).
\]
Taking its imaginary part and bounding the weighted norm gives the whole-space/box error
\[
 \frac{2}{e^{\sqrt2R}-1}\int_0^{54}V(r)e^{\sqrt2r}\dd r
 <6.162\times10^{-31}.
\]
The exact potential is linearly interpolated on $65536$ intervals. Polynomial bounds on $V''$ give interpolation error below $0.000437020$. Exact cosine moments of this interpolant are evaluated by a ball-arithmetic radix-two FFT and the piecewise-linear sinc-squared formula.

Use the first $640$ sine modes. Truncate the \emph{positive resolvent series before multiplication by $\sqrt V$}; its remainder norm is at most
$\norm V_\infty/[1+((641\pi/64)^4)]<0.000486654$. This avoids dropping an uncontrolled off-diagonal block of a preconditioned potential matrix. A floating eigensystem is only a frozen witness: ball bounds on its Gram error and eigensystem residual certify a finite-matrix spectral error below $3.227\times10^{-14}$.

Combining all errors proves
\[
 \dist\bigl(1,\spec(A^{-1/2}VA^{-1/2})\bigr)>0.04230627284>1/25.
\]
The Newton map with radial inverse norm $<25$ is invariant and contracting on $\norm{h}_Z\le4\times10^{-6}$, since
\[
 25(1.3\times10^{-7})+\frac{25\cdot211}{8}(4\times10^{-6})^2
 <4\times10^{-6}
\]
and its Lipschitz bound is below $0.006$. This proves exact radial existence.

\subsection{All angular channels, with no spectral truncation assumption}
The real part of $X$ has even angular degree and the imaginary part odd degree. For an angular harmonic $Y_\ell$, write
$\phi=r^\ell Y_\ell f(r)$ and $H_0=r^\ell Y_\ell h(r)$. The reduced radial Laplacian is $D_Df=f''+(D-1)f'/r$, $D=2\ell+3$. For each exact rational $\mu>0$, set
\[
 d_1=\Delta\phi-H_0,\qquad
 d_2=\Delta H_0+\phi-\mu^{-1}\beta|q|^{8/3}\phi.
\]
Then $K=A^{-1}\beta|q|^{8/3}$, self-adjoint in $Z$, satisfies
\[
 \norm{K\phi-\mu\phi}_Z
 \le\mu(\norm{d_1}_2^2+\norm{d_2}_2^2)^{1/2}.
\]
All radial factors have nonnegative integer powers at the origin. Noninteger powers of $q$ are enclosed by zero-safe Taylor estimates. Norm lower bounds and endpoint matching are verified from the frozen polynomials.

Each of the disjoint intervals $[\mu_j-10^{-3},\mu_j+10^{-3}]$ contains a genuine eigenvalue. An independent whole-space Green-kernel Hilbert--Schmidt bound $\tau$ gives
\[
 \mu_{\rm omitted}\le
 \left(\tau-\sum_j(\mu_j-10^{-3})^2\right)^{1/2}.
\]
Thus the residual certificates identify actual spectrum and the remaining square budget controls \emph{all} omitted spectrum. No orthogonality of numerical trials is assumed.

\begin{center}
\small
\begin{tabular}{@{}lll@{}}\toprule
Sector & Certified reference-profile information & Completion \\\midrule
Real even $\ell=0$ & radial gap $>0.04230$ & exact-profile transfer \\
Real even $\ell=2,4,6,8,10,12,14$ & gap $>0.0645$ in each sector & mixed residuals + HS \\
Real even $\ell\ge16$ & $\mathrm{HS}^2<0.702$ at $\ell=16$ & angular monotonicity \\
Imaginary odd $\ell=1$ & gap $>0.006$ & four BS intervals + HS \\
Imaginary odd $\ell=3$ & gap $>0.041$ & mixed residuals + HS \\
Imaginary odd $\ell\ge5$ & strict positivity at $\ell=5$ & angular monotonicity \\\bottomrule
\end{tabular}
\end{center}
The potential transfer is at most
$(8/3)(6.3+10^{-6})^{5/3}10^{-6}$ for $L_-$, and $11/3$ times this for $L_+$. Every exact-profile channel retains gap $>1/200$. Monotonicity is proved on the effective-dimensional radial forms; it is used only after a strictly positive channel, not across a negative channel. Taking the Hilbert direct sum proves the full inverse in Theorem~\ref{thm:n3}.

\subsection{The inverse pairing and the traveling branch}
An independently frozen dipole $Z_0=x_1f(r)$ satisfies
\[
 \norm{Z_0}_Z<808.262,\quad
 \norm{L_-(q)Z_0-\partial_1q}_{Z^*}<0.000182258,
 \quad\ip{\partial_1q}{Z_0}<-1323.3970.
\]
For $g=\partial_1Q$, $L=L_-(Q)$, $r=g-LZ_0$, the exact identity
\begin{equation}\label{eq:quadraticresponse}
 \ip g{L^{-1}g}
 =2\ip g{Z_0}-\ip{LZ_0}{Z_0}+\ip r{L^{-1}r}
\end{equation}
makes the inverse error quadratic in the response residual. The certified potential transfer gives $\norm r_{Z^*}<0.046504$ and $\norm{L^{-1}}<200$. Substitution, including the $Q-q$ derivative error, gives
$I<-1285.375835741006$, hence~\eqref{eq:n3I}.

The map $(v,U)\mapsto AU+\ii v\partial_1U-|U|^{8/3}U$ is $C^3$ from $\R\times X$ to $X^*$. The implicit-function theorem gives the stated expansion. The following direct contraction supplies an explicit interval. Write $z=L_-^{-1}\partial_1Q$, so that $\norm z_Z<818$, and let $N(w)=g(Q+w)-g(Q)-Dg(Q)w$. On $\norm w_Z\le1000|v|$, $|v|\le10^{-9}$, the real two-component derivative obeys
\[
 \norm{D^2g(\zeta)}\le\frac{88}{9}|\zeta|^{5/3}<211.
\]
The embedding and $\norm{\partial_1}_{Z\to Z^*}<3/5$ imply
\[
 \norm{N(w)}_{Z^*}\le\frac{211}{8}\norm w_Z^2,\qquad
 \operatorname{Lip}(N)\le\frac{211}{4}\,1000|v|.
\]
Because $L=L_+\oplus L_-$ is real linear, the correct fixed-point equation keeps multiplication by $\ii$ inside its inverse:
\[
 w=-\ii vz+L^{-1}N(w)-vL^{-1}(\ii\partial_1w).
\]
Its image has norm at most $818|v|+5{,}275{,}120{,}000v^2<1000|v|$, and its Lipschitz constant is at most $0.01055012$. Moreover
\[
 \norm{w+\ii vz}_Z\le5{,}275{,}120{,}000v^2,\qquad
 |P_1(U_v)-2Iv|<1.1\times10^{12}v^2.
\]
The latter follows from $\norm{\nabla f}_2\le\norm f_Z/\sqrt2$, $\norm Q_Z<137$ and the quadratic momentum expansion. The cutoff dilation proof of~\eqref{eq:pohozaev} applies to this elliptic $H^2$ profile: elliptic regularity gives $H^4$, and all cutoff errors vanish by integrable tails. Therefore
\[
 8E(U_v)\le-2560v^2+1.1\times10^{12}|v|^3\le-1460v^2,
\]
proving Theorem~\ref{thm:n3}. The noninteger nonlinearity is not claimed analytic. In particular this result does not automatically transfer the cubic $H^{10}$ gluing construction to dimension three.

The small speed range is at frequency one. Exact scaling gives
\[
 u_{\lambda,v}(t,x)=\lambda^{3/2}e^{\ii\lambda^4t}
 U_v(\lambda x-v\lambda^4te_1),
\]
with physical speed $\lambda^3v$ and energy $\lambda^4E(U_v)$. Hence any prescribed nonzero physical velocity, or any prescribed negative energy, can separately be realized. Independent simultaneous prescription of mass, energy and velocity is not asserted.

\section{A sharper four-dimensional profile certificate}
The original frozen four-dimensional profile is unchanged. Introduce a rational auxiliary $z$ on its existing mesh by nodal jets
\[
 z=q''+3q'/r,\quad z'=q'''+3q''/r-3q'/r^2,
 \quad z''=q^3-q-3z'/r,
\]
and the differentiated formula for $z'''$. At zero use the smooth radial values, and at the compact endpoint use zero jets. Septic Hermite interpolation gives an admissible global auxiliary field; these are definitions of $z$, not assertions that $q$ solves the ODE. With
\[
 d_1=\Delta q-z,\quad d_2=\Delta z+q-q^3,
 \qquad Aq-q^3=\Delta d_1+d_2,
\]
exact polynomial integration in $256$-bit balls yields
\[
 \norm{d_1}_2<1.9206685514\times10^{-8},\quad
 \norm{d_2}_2<4.1805837651\times10^{-8},\quad
 \norm{Aq-q^3}_{Z^*}<4.601\times10^{-8}.
\]
The original radial Ritz lower bounds $\mu_1\ge5.395$, $\mu_2>2.999$ and full HS-square bound $<38.80$ imply $\mu_3<0.84$, hence inverse norm $<25/4$. With $\rho=3\times10^{-7}$,
\[
 \frac{25}{4}\left(4.601\times10^{-8}+
          \frac{(312+\rho)\rho^2}{256}\right)<\rho,
\]
and the contraction constant is below $10^{-5}$. Uniqueness on the old $0.006$ ball identifies the same exact $Q$, now with
\begin{equation}\label{eq:rho4}
 \norm{Q-q}_Z<3\times10^{-7},\qquad
 \norm{A^{-1/2}(Q^2-q^2)A^{-1/2}}<2.438\times10^{-7}.
\end{equation}
This strengthens subsequent operator transfers by a factor of $20000$ in profile radius without fitting a new numerical profile.

\section{A certified radial quartet and the exact unstable dimension}\label{sec:quartet}
This additional result is independent of the new gluing mechanism. At the exact four-dimensional stationary profile, the radial Hamiltonian generator
\[
 \mathcal A_Q=\begin{pmatrix}0&L_-\\-L_+&0\end{pmatrix}
\]
has exactly one eigenvalue, counted with algebraic multiplicity, in
\begin{equation}\label{eq:rectangle}
 1.5<\Rea\lambda<4.5,\qquad4.5<\Ima\lambda<7.5.
\end{equation}
It is therefore algebraically simple in this sector, and the Hamiltonian symmetries give a simple quartet. Its numerical approximation was used to choose test functions only; the certified location is the whole rectangle~\eqref{eq:rectangle}.

\subsection{Finite block and infinite positive complement}
Quotient the phase kernel $(0,Q)$ using the $Z$-orthogonal complement. Let $\widehat A$ be the induced group generator and $H$ the bounded energy form operator there. The previous radial Morse certificate gives exactly three negative directions and invertibility of $H$, with $\norm H<6$. Choose seven smooth compactly supported trial vectors: four real quartet directions and three critical zero-chain directions. Freeze their endpoint jets and a change of basis as exact dyadic rationals. Put $W:\R^7\to D(\widehat A)$ and let $B_0$ be the explicit conjugate of two real rotation-growth blocks and a length-three nilpotent block.

The rigorous integrations certify that $G=W^*HW$ has all three negative directions. Its energy-orthogonal complement $C$ is positive and coercive. To see coercivity, let $m=\norm{H^{-1}}^{-1}$ and $P_C=I-WG^{-1}W^*H$. Extending a functional on $C$ by $P_C$ and applying $H^{-1}$ gives the compressed inverse bound $\norm{H_C^{-1}}\le\norm{P_C}/m$.

The off-diagonal blocks of $\widehat A$ in $\operatorname{Ran}W\oplus C$ are bounded: use $W^*H\widehat A=-(\widehat AW)^*H$ and $W\R^7\subset D(\widehat A)$. Removing them preserves the group-generator property. The complementary group conserves its positive energy, so its resolvent in that norm is at most $1/\Rea\lambda$.

Write $g\ge\norm{G^{-1}}$, $w\ge\norm W$, $M\ge\norm H$, $q_c=1+w^2gM$, and $\epsilon\ge\norm{\widehat AW-WB_0}$. If $B$ is the exact compressed finite matrix, the Schur correction has norm at most
\[
 gM q_c^2\epsilon^2/\Rea\lambda.
\]
Thus on a contour $\Gamma$ with $\Rea\lambda\ge\alpha_\Gamma>0$, the condition
\begin{equation}\label{eq:rouche}
 k\left(\delta_B+\frac{gM q_c^2\epsilon^2}{\alpha_\Gamma}\right)<1,
 \quad k\ge\sup_\Gamma\norm{(B_0-\lambda)^{-1}},\quad
 \delta_B\ge\norm{B-B_0},
\end{equation}
preserves the finite Schur determinant's zero count. Analytic Schur factorization transfers algebraic multiplicity to the full generator. This controls the entire infinite complement without discretizing it.

\subsection{Exact-profile transfer and the verified margin}
The frozen trial functions are $C^5$ compact radial polynomials of degree eleven. Origin odd jets through order five and endpoint jets through order five are checked exactly. At $q$, exact polynomial integrations give the Gram inertia, the signed compressed matrix, and the whole-space $H^2$ strong residual. In particular $\norm{B_q-B_0}<9.706\times10^{-9}$; estimating this quantity only by the strong residual would lose decisive cancellation.

For $a=Q-q$,~\eqref{eq:rho4}, the original strong residual $\norm{Aq-q^3}_2<0.000242$, and a Fourier interpolation estimate give
\[
 \norm{Aa}_2<0.0003,\qquad\norm a_\infty<8.4\times10^{-7}.
\]
Explicit product estimates for $(Q^2-q^2)U=(2qU)a+Ua^2$ transfer every strong residual to the exact profile; $\norm{\nabla a}_4^2\le\norm a_Z\norm{Aa}_2/16$ controls the quadratic derivative term. The phase quotient error is bounded by $\norm{Q-q}_Z$ times the corresponding fixed coordinate row norm. Gram and compressed-matrix transfers give the following conservative bounds:
\begin{center}
\begin{tabular}{@{}lc@{}}\toprule
Quantity & Verified upper bound \\\midrule
$\epsilon$ & $0.002991362$ \\
$g$ & $5.858206$ \\
$w$ & $1.000001876$ \\
$q_c$ & $36.149366$ \\
$\delta_B$ & $0.000877055$ \\
$k$ & $1.330629$ \\\bottomrule
\end{tabular}
\end{center}
Each contour side is covered by $160$ closed interval segments, with rigorous complex matrix inversion over each entire segment. For $M=6$ and $\alpha_\Gamma=1.5$, the left side of~\eqref{eq:rouche} is less than $0.365771$. The reference block has exactly one simple eigenvalue inside; therefore so does the continuum radial generator.

These two real unstable dimensions persist under sufficiently small traveling speed. Together with the previously established $29$ nonradial unstable dimensions and the moving-wave Hamiltonian index upper bound $31$, they give
\[
 \dim_{\R}E^u(Q_v)=\dim_{\R}E^s(Q_v)=31
\]
for sufficiently small nonzero $v$. This certification does not give an explicit velocity threshold for that exact dimension statement.

\section{Time-averaged kinetic control and intrinsic separation}
This section is independent of every profile and interval certificate. Let a forward-global strong $H^2$ solution of~\eqref{eq:pde} have $M>0$, $E<0$, $P=0$. Put
\[
 K(t)=\norm{\Delta u(t)}_2,\quad h(t)=1+K(t)^2,\quad
 \Lambda=16M|E|,\quad\theta=\min(4/N,1/4),
\]
\[
 d_R(t)=M^{-2}\iint_{|x-y|>R}|u(t,x)|^2|u(t,y)|^2\dd x\dd y.
\]

\subsection{A finite-window virial inequality}
Choose a smooth radial $\phi_R$ with quadratic core, bounded gradient $O(R)$, constant exterior and $D^2\phi_R\le I$. Set
\[
 I_R=2\iint|u(y)|^2\nabla\phi_R(x-y)\cdot\Ima(\bar u(x)\nabla u(x))\dd x\dd y.
\]
The localized fourth-order identity of~\cite{BL} and its interaction version, as used in~\cite{Komada}, give the following bounds with the kinetic dependence retained:
\begin{equation}\label{eq:instant}
 |I_R|\le ARh^{1/4},\qquad
 I_R'\le-\Lambda+BR^{-2}h^{1/2}+Chd_R^\theta.
\end{equation}
Here $A,B,C$ depend only on dimension, mass and cutoffs, not on a temporal supremum of $K$.

For clarity, the closure uses mass tails alone. With $G=M^{1/4}K^{1/2}$ and $H_R=K+cG/R+c\sqrt M/R^2$, localized Gagliardo--Nirenberg bounds the nonlinear tail by $C_MH_R^2d_R^{\min(4/N,1)}$. The cross-flux tail is bounded by $C_MK^{3/2}H_R^{1/2}d_R^{1/4}$. These are at most $Chd_R^\theta$. The leading Hessian term has a favorable upper sign, so no smallness of curvature tails is needed. The derivative falling on the second mass density is retained; its uncut contribution is $-8P\cdot W$, $W=\Ima\int\Delta u\nabla\bar u$, and vanishes here. Remaining cutoff errors are $BR^{-2}h^{1/2}$. Smoothing on compact time intervals justifies the integrated identity in $H^2$.

On $J=[S,S+T]$ define $H_J=T^{-1}\int_Jh$. Choose $a$ in the first quarter and $b$ in the last quarter with $h(a),h(b)\le4H_J$. Then $b-a\ge T/2$. Integrating~\eqref{eq:instant} and using Jensen gives
\begin{equation}\label{eq:finitewindow}
 \boxed{\quad\frac\Lambda2\le
 A_0\frac{RH_J^{1/4}}T+B\frac{H_J^{1/2}}{R^2}
 +\frac CT\int_Jh(t)d_R(t)^\theta\dd t.\quad}
\end{equation}
This holds on every finite window inside the lifespan, without a global kinetic hypothesis.

\subsection{Mass transport closes the first-moment endpoint}
Choose smooth $\psi$ equal to zero on $|z|\le1/2$ and one on $|z|\ge1$, with $0\le\psi\le1$. The smoothed pair tail
\[
 F_R=M^{-2}\iint\psi((x-y)/R)|u(x)|^2|u(y)|^2\dd x\dd y
\]
satisfies $d_R\le F_R\le d_{R/2}$. Testing the density equation with real $f$ gives
\[
 \int f\partial_t|u|^2
 =2\Ima\int(\Delta f)\bar u\Delta u
       +4\Ima\int\nabla f\cdot(\Delta u\nabla\bar u).
\]
It follows that $|F_R'|\le D(R^{-1}K^{3/2}+R^{-2}K)$ and hence
\begin{equation}\label{eq:transport}
 |F_R(t)-F_R(s)|\le D\left[
 R^{-1}\Big(\int_JK^2\Big)^{3/4}|t-s|^{1/4}
 +R^{-2}\Big(\int_JK^2\Big)^{1/2}|t-s|^{1/2}\right].
\end{equation}
The flux exponent $3/2$ is strictly below the kinetic budget exponent $2$.

\begin{theorem}[Separation on an entire time interval]\label{thm:average}
Suppose for some $H_0,T_*<\infty$,
\[
 \sup_{S\ge0,\,T\ge T_*}\frac1T\int_S^{S+T}(1+K(t)^2)\dd t\le H_0.
\]
There are $c,\delta,\eta,T_0>0$ such that every window $[S,S+T]$, $T\ge T_0$, contains an interval of length at least $\eta T$ on which, simultaneously for every $a\in\R^N$,
\begin{equation}\label{eq:allcenters}
 \int_{|x-a|>cT}|u(t,x)|^2\dd x\ge\delta M.
\end{equation}
No spatial moment, radiality, higher kinetic time moment or uniform integrability is assumed.
\end{theorem}
\begin{proof}
Set $R=c_0T$ with $c_0$ small. The first two errors in~\eqref{eq:finitewindow} are absorbed for large $T$, giving a positive lower bound on $T^{-1}\int_Jhd_R^\theta$. Since $H_J\le H_0$, some $t_0\in J$ satisfies $d_R(t_0)\ge\delta_0>0$. For $|t-t_0|\le\eta T$,~\eqref{eq:transport} is bounded by
\[
 D(c_0^{-1}H_0^{3/4}\eta^{1/4}
       +c_0^{-2}H_0^{1/2}T^{-1}\eta^{1/2}).
\]
Choose $\eta$ small to make this at most $\delta_0/2$. Then $d_{R/2}\ge F_R\ge\delta_0/2$ on an interval of length $\eta T$ within $J$. Finally
$d_{2r}\le(2/M)\int_{|x-a|>r}|u|^2$ for every center $a$, which proves~\eqref{eq:allcenters} after renaming constants.
\end{proof}

There is also an unconditional finite-window alternative: either $H_J>(T/L)^2$, or an interval of length at least
$c_\eta H_J^{-4(1+1/\theta)}T$ satisfies
\[
 d_{(a_0/2)TH_J^{-1/4}}\ge c_\delta H_J^{-1/\theta}.
\]
This follows by taking $R=a_0TH_J^{-1/4}$ in~\eqref{eq:finitewindow} and then applying~\eqref{eq:transport}; all constants are independent of the window.

\subsection{Localization forces averaged grow-up}
\begin{theorem}[Uniform long-window grow-up]\label{thm:grow}
If the same global solution is uniformly tight in mass modulo translations, namely
\[
 \lim_{R\to\infty}\sup_{t\ge0}\inf_a
 \frac1M\int_{|x-a|>R}|u(t,x)|^2\dd x=0,
\]
then
\[
 \lim_{T\to\infty}\inf_{S\ge0}\frac1T\int_S^{S+T}K(t)^2\dd t=\infty.
\]
If the mass tail is uniformly at most $DR^{-q}$, $0<q\le2$, then, for all sufficiently large $T$ and every $S\ge0$,
\begin{equation}\label{eq:rate}
 \frac1T\int_S^{S+T}K(t)^2\dd t
 \ge cT^{4q\theta/(4+q\theta)}.
\end{equation}
\end{theorem}
\begin{proof}
Mass tightness implies $\mu(R):=\sup_td_R(t)\to0$. If arbitrarily long windows had bounded $H_J$, take $R=cT$ in~\eqref{eq:finitewindow}. The endpoint term is small for small $c$, the cutoff term tends to zero, and the tail term is at most $CH_J\mu(cT)^\theta\to0$, a contradiction.

For a power tail let $\alpha=q\theta\le1/2$ and take $R=L_1H_J^{1/\alpha}$. A sufficiently large fixed $L_1$ absorbs both the cutoff and tail terms, since $H_J^{1/2-2/\alpha}\le1$ and $H_JR^{-\alpha}=L_1^{-\alpha}$. The remaining inequality is $\Lambda/4\le C H_J^{1/4+1/\alpha}/T$. Rearrangement and subtraction of the one in $H_J$ prove~\eqref{eq:rate}.
\end{proof}
The conclusion is about time averages, not pointwise convergence of $K(t)$ to infinity. Algebraic tails with $q\le2$ need not imply finite variance.

\section{Relations to prior work and boundaries of the conclusions}\label{sec:context}
The construction combines existing Hamiltonian structure, variational identities and rigorous spectral approximation. Lin--Zeng~\cite{LZ} supplies the triangular decomposition; the deduction here is the compact finite-rank metric with an arbitrarily small center growth rate and its moving-core use. Excited multi-soliton construction has substantial precedent, including the high-speed NLS work of C\^ote--Le Coz~\cite{CL}. That result is not itself an arbitrary-small-speed gluing theorem for these fourth-order negative-energy profiles.

The interaction virial also has precedent. Our pure analytic extension changes the closure: low-kinetic endpoint selection and the mass-transport time modulus require only first time averages and yield intervals of separation. The comparison with Komada's compact-orbit argument does not amount to a claim of priority for the interaction functional itself.

There is a material published conflict. Gou's mass-critical cylindrical statement~\cite[Theorem 1.2]{Gou}, at zero lower-order dispersion, asserts finite-time blowup or infinite-time $H^2$ grow-up for negative-energy data in a cylindrical class with finite axial moment. Theorem~\ref{thm:pair}, if verified as presented, lies in that class and is incompatible with that conclusion. The relevant virial weight is $\phi_R(y,s)=\psi_R(|y|)+s^2/2$. Its axial gradient is unbounded. The usual bounded-gradient estimate cannot discard
\[
 2\norm{su}_2\norm{\partial_su}_2.
\]
An initially finite axial moment does not bound this factor uniformly in time. In our pair, the variance grows quadratically while the $H^2$ norm stays bounded. The companion audit identifies this closure issue explicitly. A proof gap by itself is not a counterexample; the independent existence and certificate arguments remain necessary.

The present work does not establish an open basin of negative-energy global data, a classification of all negative-energy dynamics, the same traveling-wave construction in every dimension, or a smooth cubic-type stable-manifold theory for the noninteger three-dimensional nonlinearity. Arbitrary higher-order dispersions and general nonlinearities remain possible extensions, not proved results of this manuscript.

\section{Reproducibility and dependency map}
The accompanying package contains the frozen inputs, interval verifiers, analytic proof notes and prior companion inputs required for the four-dimensional application. The following separation is maintained:
\begin{center}
\begin{tabular}{@{}>{\raggedright\arraybackslash}p{0.32\textwidth}>{\raggedright\arraybackslash}p{0.60\textwidth}@{}}\toprule
Conclusion & Necessary input \\\midrule
Three-dimensional traveling wave & exact profile residual; radial inverse; all even/odd angular channels; inverse pairing; analytic IFT and Pohozaev identity \\
Four-dimensional zero-momentum pair & previously certified negative-energy profile; compact metric lemma; moving-center estimates; weighted nonlinear contraction \\
Improved four-dimensional error radius & unchanged rational profile; mixed residual; original radial Ritz and HS estimates \\
Averaged spreading and grow-up & strong $H^2$ flow; conservation; interaction virial; mass transport; no numerical data \\\bottomrule
\end{tabular}
\end{center}
Rational polynomial data and binary floating witnesses are interpreted exactly before interval validation. Numerical eigenvalues are never accepted as a proof of an infinite-dimensional gap. Whole-space tails, origin regularity, compact-support matching, omitted spectral mass and transfer from $q$ to $Q$ are separate checks. A successful script verifies its stated inequalities, not every analytic assertion in the manuscript. Independent external replication of both parts remains a distinct step.

\begin{thebibliography}{9}
\small\setlength{\itemsep}{2pt}\setlength{\parskip}{0pt}
\bibitem{LZ} Z. Lin and C. Zeng, \emph{Instability, index theorem, and exponential trichotomy for linear Hamiltonian PDEs}, Memoirs of the AMS 275 (2022), no. 1347. \url{https://arxiv.org/abs/1703.04016}.
\bibitem{BL} T. Boulenger and E. Lenzmann, \emph{Blowup for biharmonic nonlinear Schr\"odinger equations}. \url{https://arxiv.org/abs/1503.01741}.
\bibitem{Komada} K. Komada, \emph{Scattering of group-invariant solutions below the ground state for the fourth-order NLS}. \url{https://arxiv.org/abs/2609.33381}.
\bibitem{CL} R. C\^ote and S. Le Coz, \emph{High-speed excited multi-solitons in nonlinear Schr\"odinger equations}, J. Math. Pures Appl. 96 (2011), 135--166. \url{https://arxiv.org/abs/1007.3034}.
\bibitem{Gou} T. Gou, \emph{Blowup of cylindrically symmetric solutions for biharmonic NLS}, arXiv:2205.08167v3. \url{https://arxiv.org/abs/2205.08167}.
\end{thebibliography}
\end{document}
